Diese Arbeit befasst sich mit der kurzfristigen Wasserstandsvorhersage für die Donau in Korneuburg, Österreich. Ziel ist es, verschiedene Vorhersagemodelle für die nächsten 24 Stunden zu vergleichen und den Einfluss flussaufwärts gelegener Messstationen sowie von Wetterdaten auf die Vorhersagegenauigkeit zu untersuchen. Die vorhergesagten Werte könnten von FloodAlert genutzt werden, um die zukünftige Wasserstandsentwicklung auf dessen Website darzustellen.

Die Wasserstands- und Wetterdaten werden vorverarbeitet und um abgeleitete Merkmale wie zeitverzögerte Werte, Wasserstandsänderungen und rollende Statistiken ergänzt. Lineare Modelle, baumbasierte Modelle und neuronale Netze werden trainiert und mit einem Persistenzmodell als Referenz verglichen. Die Modellauswahl erfolgt anhand einer \ac{CV} mit einem sich schrittweise erweiternden Trainingsfenster, während ein chronologisch abgetrennter Testdatensatz ausschließlich für die abschließende Bewertung verwendet wird.

Insgesamt erzielt XGBoost mit $12{,}43$ cm den niedrigsten \ac{RMSE} in der \ac{CV} und wird daher als finales Modell ausgewählt. Auf dem Testdatensatz erreicht XGBoost einen \ac{RMSE} von $13{,}90$ cm, während das \ac{MLP} mit einem \ac{RMSE} von $12{,}11$ cm besser abschneidet. Die Einbeziehung flussaufwärts gelegener Messstationen und von Wetterdaten reduziert den \ac{RMSE} in der \ac{CV} um $25{,}8\%$ gegenüber der ausschließlichen Verwendung von Merkmalen der Zielstation ohne Wetterdaten. Das Persistenzmodell liefert für die ersten sechs Stunden bessere Ergebnisse, während XGBoost über den verbleibenden Vorhersagehorizont besser abschneidet. Allerdings hat XGBoost Schwierigkeiten, extreme Wasserstände genau vorherzusagen. Oberhalb der Alarmschwelle von 545 cm steigt der \ac{RMSE} auf $71{,}44$ cm, wobei das Modell den Wasserstand systematisch unterschätzt.

Abschließend beschreibt ein Konzept für den produktiven Einsatz, wie das Modell in FloodAlert integriert werden könnte. Die Vorhersagen könnten zusätzliche Informationen über zukünftige Wasserstände liefern. Die Vorhersagegenauigkeit bei Extremereignissen reicht jedoch nicht aus, um die Vorhersagen als alleinige Grundlage für Hochwasserwarnungen zu nutzen. Darüber hinaus sind die Ergebnisse auf eine einzelne Messstation beschränkt.
